\section{Falsification and Reporting Rules}

The richer model hierarchy is not supported merely because PI or PID occasionally has the lowest RMSE. Generating-family recovery concerns frozen validation choices on independent synthetic replicates; final-test predictive advantage is a separate outcome. High family recovery can coexist with practically small loss differences. Conversely, low recovery need not imply poor prediction. Neither a nominal winner nor predictive equivalence alone establishes structural identifiability.

The study reports weak family recovery, simpler-baseline selection, uncertain parameter estimation and practical predictive equivalence separately. A chance benchmark requires a stated candidate-eligibility and random-selection rule. Rank failure, insufficient observations and parameter degeneracy remain explicit design/numerical outcomes. These findings restrict interpretation in the investigated generator universe rather than refuting a human mechanism.

The manuscript will not map P, PI, or PID outcomes onto attachment categories or infer human mechanisms from synthetic recovery. Human scale validation and longitudinal observations belong exclusively to M03 and M04.
